\begin{proposition}[An improved bound in degree zero]
\label{c:improved-zero-weight}
\status{Proved; optimality for general coefficients of weight at least three remains unresolved}
Let $f:X\to Y$ be separated and of finite type over $\Z[1/\ell]$
or over a finite field of characteristic different from $\ell$, with
$X$ regular, and let $y\in Y$ have finite residue field.  Let $F$ be
strictly special of weight $j\geq0$, in either convention of
Definition~\ref{c:level-def}.  Put
\[
 b(0)=0,\qquad b(1)=1,\qquad b(j)=2j-2\quad(j\geq2).
\]
Then the eigenvalues of geometric Frobenius on
$(f_*F)_{\bar y}$ have weights at most $b(j)$.
In particular, the expected bound $j$ holds for $j\leq2$.
The assertion about weights does not replace the geometric level
definition or assert a corresponding smaller level.
\end{proposition}

\begin{proof}
We first prove the local-field input, including the two extreme
monodromy maps.  Let $K$ be a finite extension of $\Q_p$ or of
$\F_p((t))$, and let $Z/K$ be smooth projective of dimension $d$.
After a finite extension and a projective semistable alteration,
the pullback on cohomology is split injective by trace divided by
the generic degree.  Thus we may use a projective strictly
semistable model.  This reduction is
\citet[Theorem~6.5, pp.~83--86]{c_dejong1996}; the assertion
descends to the original inertia invariants because Frobenius over
the extension acts there as a power of the original Frobenius.

Let $Y^{(r)}$ be the disjoint union of the $(r+1)$-fold intersections
of the special-fibre components.  We use the weight spectral sequence
\[
 E_1^{a,b}=\bigoplus_{r\geq\max(0,-a)}
 H^{b-2r}(Y^{(a+2r)}_{\bar k},\Q_\ell)(-r)
 \Longrightarrow H^{a+b}(Z_{\bar K},\Q_\ell).
\]
The $E_1^{a,b}$ are pure of weight $b$, the sequence degenerates at
$E_2$, and $N$ maps $(a,b)$ to $(a+2,b-2)$, with the matching
Tate twist.  These are
\citet[Satz~2.10 and (2.11), pp.~40--41]{c_rapoportzink1982}.
For $H^d$, all weights lie in $[0,2d]$.

The extreme map
\[
 N^d:E_2^{-d,2d}\longrightarrow E_2^{d,0}(-d)
\]
is an isomorphism.  Suppressing twists, it is the natural map
\[
 \ker\bigl(H^0(Y^{(d)})\xrightarrow{\mathrm{Gys}}
 H^2(Y^{(d-1)})\bigr)
 \longrightarrow
 \operatorname{coker}\bigl(H^0(Y^{(d-1)})
 \xrightarrow{\mathrm{Res}}H^0(Y^{(d)})\bigr).
\]
Restriction and Gysin are adjoint.  On the rational vector space
of functions on the finite set $Y^{(d)}$, the pairing is a positive
sum of products.  Its restriction to the rational image of
restriction is nondegenerate, and remains so over $\Q_\ell$.
The displayed kernel is its orthogonal complement and maps
isomorphically to the quotient.  This is the argument of
\citet[Proposition~1.5, p.~9]{c_yamamura2026weights}.

When $d\geq2$, the next extreme map
\[
 N^{d-1}:E_2^{-d+1,2d-1}\longrightarrow
 E_2^{d-1,1}(-(d-1))
\]
is also an isomorphism.  Put $U=H^1(Y^{(d-1)})$, so the components
in this expression are smooth projective curves, and put
$R=\operatorname{im}[H^1(Y^{(d-2)})\to U]$.
The signed restriction map is induced by the corresponding
homomorphism from the product of the reduced Picard varieties of
the surface strata to the product of the curve Jacobians.
Consequently $R$, with the usual Tate twist understood, is the
Tate realization of an abelian subvariety of that product.
The product principal polarization restricts to a polarization on
this abelian subvariety.  Thus the cup pairing on $U$ is
nondegenerate on $R$.  Since Gysin is the adjoint of restriction,
its kernel is $R^\perp$, and $R^\perp\to U/R$ is an isomorphism.
This proves the next extreme map, with all signs included.
We have used the Picard-polarization argument in the proof of
\citet[Satz~2.13, pp.~41--43, especially p.~42]{c_rapoportzink1982},
not an application of its surface theorem to a higher-dimensional
variety.  The resulting general-dimensional statement is
\citet[Proposition~1.7, p.~10]{c_yamamura2026weights}.

Inertia is unipotent in this model, so its invariants are $\ker N$.
An element in the kernel cannot have a nonzero weight-$2d$ component,
and for $d\geq2$ it cannot have a weight-$(2d-1)$ component either:
on each of these components a positive power of $N$ is injective.
The weight filtration is strict here, since it is the decomposition
into disjoint generalized Frobenius eigenspaces by their weights.
Hence $H^d(Z)^{I_K}$ has weights at most $2d-2$ for $d\geq2$, and
at most $1$ for $d=1$.  Degree zero has weight zero.  This recovers
the middle-degree consequences of
\citet[Theorem~1.6 and Proposition~1.8, pp.~9--11]{c_yamamura2026weights}.
The proof just given works in both local-field characteristics.
For a direct summand of $H^j$ of a smooth projective variety of
dimension at least $j$, the cut-and-pairing construction of
Lemma~\ref{c:proper-level} realizes it as a direct summand in degree
$j$ of a smooth projective $j$-fold, after a finite extension.
If the dimension is $d<j$, replace $Z$ by
$Z\times\mathbb P^{j-d}$: the $H^0$ projector on the projective
factor splits off the original $H^j$, and the product has dimension
$j$.  When $j>2d$ the original group is already zero.
Thus the same bound $b(j)$ applies.
We use only these independently proved middle-degree consequences.
The displayed higher-degree hard-Lefschetz isomorphism on p.~10
of the cited preprint prints the twist $i-d$; in our convention
the correct twist is $d-i$.  Its subsequent numerical weight bounds
use the latter sign.  This printed slip is not used in our proof.

We now reduce the degree-zero stalk to these local fields.
Write $S=\Spec\mathcal O_{Y,\bar y}^{\mathrm{sh}}$ and
\[
 V=(f_*F)_{\bar y}=H^0(X\times_Y S,F).
\]
There are finitely many connected components; they are integral
because the source is regular.  Descend these finitely many
open-and-closed components and a compactification to a finite
\'etale neighbourhood of $y$.  We do not assert that every
$\ell$-adic section descends to one such neighbourhood.
After a finite residue extension the relevant components are fixed.
The closure of each component in a proper compactification meets
the closed fibre.  Choose a closed point $c$ there.  A component
lying over $y$ can be evaluated at a finite closed point directly,
giving only weight $j$.

For the other components we construct a curve approaching $c$ from
the chosen dense open.  In the local domain $A$ of its closure at
$c$, choose $0\ne g\in\mathfrak m_A$ with $D(g)$ contained in
that open; in the horizontal mixed-characteristic case include $p$
as a factor of $g$.  Since $A$ is an excellent local domain,
the equidimensional dimension formula gives
$\dim A/(g)=\dim A-1$; thus $g$ extends to a system of parameters
$g,a_1,\ldots,a_{d-1}$ of $A$.  A minimal prime $P$ of
$(a_1,\ldots,a_{d-1})$ has height at most $d-1$ and cannot
contain $g$, since that would make $P=\mathfrak m_A$.
Moreover $g$ is primary for the maximal ideal of $A/P$, so
$\dim(A/P)=1$.  The resulting integral curve has generic point
in the required open.  Normalize, localize at a point above $c$,
and complete.  Excellence gives a discrete valuation ring with
finite residue field, whose fraction field is a finite extension
of $\Q_p$ or of $\F_p((t))$.  The entire curve construction is
performed after the preceding finite descent, over a scheme of
finite type over $\Z$ or a finite field.

Let $T$ be this trait.  The map $T^{\mathrm{sh}}\to S$, and its
generic-point map to the chosen component, pull back a section in
$V$ to an inertia-invariant vector of the restricted sheaf.
This map is equivariant for the residue Galois group after the
finite residue extension.  Evaluation at one geometric point of
a connected component is injective for a lisse sheaf.  Taking
the finite sum of these evaluations therefore embeds $V$ into
the local-field invariant spaces just considered and the pure
weight-$j$ evaluation spaces.  Since $j\leq b(j)$ for every
$j\geq0$, their weights are all at most $b(j)$.
Restriction of a strictly special sheaf retains its presenting
smooth projective family and its projector, so the local-field
input applies.  Finally the bound descends through the finite
residue extension by taking powers of Frobenius.
\end{proof}

\begin{corollary}[The first remaining lower-endpoint case]
\label{c:threefold-weight-gate}
\status{Proved equivalence for this case; the stated monodromy map remains unresolved in general}
Let $Z$ be a smooth projective threefold over a mixed-characteristic
local field with a projective strictly semistable model, and
$\ell\ne p$.  The following are equivalent:
\begin{enumerate}[label=(\roman*)]
\item $H^3(Z_{\bar K},\Q_\ell)^{I_K}$ has weights at most $3$;
\item $N:\operatorname{Gr}^W_4H^3\to
       \operatorname{Gr}^W_2H^3(-1)$ is injective;
\item the weight--monodromy conjecture holds for $H^3(Z_{\bar K},\Q_\ell)$.
\end{enumerate}
Here the weight filtration $W$ is supplied by the weight spectral
sequence, and the weight--monodromy statement means that
$N^r:\operatorname{Gr}^W_{3+r}\to
\operatorname{Gr}^W_{3-r}(-r)$ is an isomorphism for every $r\geq0$.
\end{corollary}

\begin{proof}
The preceding proof already makes $N^3$ from weight $6$ to weight
$0$, and $N^2$ from weight $5$ to weight $1$, isomorphisms.
Consequently the only possible weight above $3$ in $\ker N$ is
weight $4$.  Strictness with respect to weights identifies its
weight-$4$ part with the kernel in (ii), proving (i)$\Leftrightarrow$(ii).
Poincar\'e duality on $H^3$, with values in $\Q_\ell(-3)$, pairs
weight $4$ perfectly with weight $2$.  Thus the source and target
in (ii) have equal dimension; injectivity is equivalent to an
isomorphism.  Together with the already verified $r=2,3$ cases,
and the identity for $r=0$, this is exactly (iii).
\end{proof}

